\section{The relative Picard functor and its \'etale sheafification}
\label{mk-chap:Picard_RelPicFunctor}

\section{Group structure on the relative Picard quotient}
\label{mk-sec:relpic_group}

The Picard group \(\Pic(X)\) of any scheme is canonically an abelian group under
tensor product of line bundles (Stacks tag 01CR; the inverse of \([\mathcal{L}]\)
is $[\mathcal{L}^{-1}] = [\mathcal{H}om_{\mathcal{O}_X}(\mathcal{L},
\mathcal{O}_X)]$). The pullback map
\(\pi_T^* : \Pic(T) \to \Pic(C \times_k T)\) of
\ref{mk-def:pullback_along_projection} respects this structure: pullback of
\(\mathcal{O}_X\)-modules sends \(\mathcal{O}_T\) to \(\mathcal{O}_{C \times_k T}\)
and is multiplicative on tensor products.

\begin{lemma}[Group structure descends to the relative Picard quotient]\label{mk-lem:rel_pic_sharp_groupoid}
  \dcref{01CR,custom}

  \uses{mk-def:line_bundle_on_product, mk-def:pullback_along_projection,
    mk-thm:relative_pic_quotient_well_defined}
  \textit{Source: \cite{kleiman-picard}, ``The Picard scheme'', \S 2,
  Defs.~df:aPf + df:Pfs (the Picard functors are abelian-group-valued, and the
  relative one is a quotient of abelian groups; cf.\ Stacks tag 01CR for the
  abelian-group structure on \(\Pic\)).}
  Let \(C/k\) be a smooth proper geometrically integral curve and let \(T\) be a
  \(k\)-scheme. The Picard group \(\Pic(C \times_k T)\) is an abelian group under
  tensor product of line bundles, and the pullback map
  \[
    \pi_T^* \;\colon\; \Pic(T) \;\longrightarrow\; \Pic(C \times_k T)
  \]
  of \ref{mk-def:pullback_along_projection} is a group homomorphism. Its image
  \[
    H_T \;\;:=\;\; \pi_T^* \Pic(T)
  \]
  is therefore a subgroup of \(\Pic(C \times_k T)\), and the quotient set of
  \ref{mk-thm:relative_pic_quotient_well_defined}
  \[
    \Pic^\sharp_{C/k}(T) \;\;=\;\; \Pic(C \times_k T) \, / \, H_T
  \]
  inherits a canonical abelian-group structure under which the quotient map
  \(\Pic(C \times_k T) \twoheadrightarrow \Pic^\sharp_{C/k}(T)\) is a surjective
  homomorphism with kernel exactly \(H_T\).
\end{lemma}
\begin{proof}
  \uses{mk-def:line_bundle_on_product, mk-def:pullback_along_projection,
    mk-thm:relative_pic_quotient_well_defined, mk-lem:tensorobj_lift_onproduct,
    mk-lem:tensorobj_preserves_locally_trivial, mk-lem:tensorobj_assoc_iso,
    mk-lem:tensorobj_comm_iso, mk-lem:tensorobj_unit_iso,
    mk-lem:tensorobj_isoclass_commgroup, mk-lem:tensorobj_inverse_invertible,
    mk-lem:pullback_unit_iso, mk-lem:pullback_tensor_iso_loctriv}


The construction proceeds in four steps over the locally-trivial carrier
  \[
    \mathrm{OnProduct}(\pi_C, \pi_T) \;=\;
      \bigl\{\, M : \text{a module on } C \times_k T \;\bigm|\;
      M \text{ is locally trivial} \,\bigr\}
  \]
  of \ref{mk-def:line_bundle_on_product}, on whose isomorphism classes the
  relative Picard quotient is taken. The comparison isomorphism
  \ref{mk-lem:pullback_tensor_iso_loctriv} makes \(\pi_T^*\) additive, while
  \ref{mk-lem:tensorobj_inverse_invertible} supplies closure under inverses.

  \emph{Step 1 (abelian group on the carrier).} On isomorphism classes of the
  carrier \(\mathrm{OnProduct}(\pi_C, \pi_T)\), addition is the descent of the
  scheme-level tensor product of line bundles,
  \[
    [L] + [L'] \;:=\; [\, L \otimes_{\mathcal{O}_{C \times_k T}} L' \,],
  \]
  where \(L \otimes L'\) is the locally-trivial descent of the scheme-level
  tensor product (\ref{mk-lem:tensorobj_lift_onproduct}); the sum lands again in
  the carrier because the tensor product of two locally-trivial modules is
  locally trivial (\ref{mk-lem:tensorobj_preserves_locally_trivial}). The zero
  element is the class \([\mathcal{O}_{C \times_k T}]\) of the structure sheaf,
  the monoidal unit, which is locally trivial. Associativity, commutativity,
  and left/right unitality of \(+\) descend from the scheme-level coherence
  isomorphisms --- the associator (\ref{mk-lem:tensorobj_assoc_iso}), the
  braiding (\ref{mk-lem:tensorobj_comm_iso}), and the unitors
  (\ref{mk-lem:tensorobj_unit_iso}) --- exactly as for the absolute Picard group
  on tensor-invertible isomorphism classes
  (\ref{mk-lem:tensorobj_isoclass_commgroup}). The additive inverse is
  \[
    -[L] \;:=\; [\, L^{\mathrm{inv}} \,],
  \]
  where \(L^{\mathrm{inv}}\) is the witness supplied by
  \ref{mk-lem:tensorobj_inverse_invertible}: it returns a \emph{locally-trivial}
  module together with an isomorphism
  \(L \otimes L^{\mathrm{inv}} \cong \mathcal{O}_{C \times_k T}\), so the inverse
  stays inside the loc-triv carrier and the group is closed under negation.

  \emph{Step 2 (the pullback homomorphism \(\pi_T^*\)).} Pullback along the
  projection \(\pi_T\) gives an additive map
  \[
    \pi_T^* \;\colon\; \Pic(T) \;\longrightarrow\; \mathrm{OnProduct}(\pi_C, \pi_T),
  \]
  whose underlying assignment of line bundles is
  \ref{mk-def:pullback_along_projection}. It is a homomorphism of abelian groups:
  it sends the unit to the unit, since
  \(\pi_T^*\mathcal{O}_T \cong \mathcal{O}_{C \times_k T}\) by the unconditional
  pullback--unit isomorphism (\ref{mk-lem:pullback_unit_iso}); and it is additive,
  since pullback commutes with tensor product on locally-trivial modules,
  \[
    \pi_T^*(N \otimes_{\mathcal{O}_T} N') \;\cong\;
      \pi_T^* N \otimes_{\mathcal{O}_{C \times_k T}} \pi_T^* N',
  \]
  by the locally-trivial comparison isomorphism
  \ref{mk-lem:pullback_tensor_iso_loctriv} specialised to the projection
  \(\pi_T\) on the \(\mathrm{OnProduct}\) carrier. This is the exact analogue,
  for the geometric pullback, of the additive monoid homomorphism on Picard
  groups induced by a ring map respecting tensor product, and is modelled on it
  field-for-field. Write \(H_T := \mathrm{im}(\pi_T^*)\) for its image, a
  subgroup of the carrier's isomorphism-class group.

  \emph{Step 3 (setoid reconciliation).} The set-valued quotient of
  \ref{mk-thm:relative_pic_quotient_well_defined} is taken modulo the
  isomorphism-class equivalence that declares two line bundles related when
  their classes differ by an element of \(H_T\). This relation coincides with
  the left-coset relation of the subgroup \(H_T\) inside the abelian group of
  Steps~1--2: both encode
  \[
    L \sim L' \quad\Longleftrightarrow\quad
      [\, L^{\mathrm{inv}} \otimes L' \,] \in H_T .
  \]
  Establishing this equivalence of relations is a concrete isomorphism-class
  argument over the carrier; it consumes the tensor inverse (Step~1) and the
  group law, but not the computational content of the comparison isomorphism.

  \emph{Step 4 (transport).} By Step~3 the quotient set underlying
  \(\Pic^\sharp_{C/k}(T)\) is in canonical bijection with the group quotient
  \(\mathrm{OnProduct}(\pi_C, \pi_T) / H_T\), which carries the canonical
  abelian-group structure of a quotient of an abelian group by a subgroup.
  Transporting this structure along the bijection of Step~3 endows
  \(\Pic^\sharp_{C/k}(T)\) with its abelian-group structure, under which the
  canonical surjection
  \(\Pic(C \times_k T) \twoheadrightarrow \Pic^\sharp_{C/k}(T)\) is a
  homomorphism with kernel exactly \(H_T\).
\end{proof}

\subsection{Helper constructions underlying the group law}
\label{mk-subsec:relpic_group_helpers}


\begin{definition}[Tensor product on relative line bundles]\label{mk-def:rel_tensor_obj}
  \dcref{custom}

  \uses{mk-def:line_bundle_on_product, mk-lem:tensorobj_lift_onproduct,
    mk-lem:tensorobj_preserves_locally_trivial}
  For relative line bundles \(L, L'\) on the locally-trivial carrier
  \(\mathrm{OnProduct}(\pi_C, \pi_T)\) of \ref{mk-def:line_bundle_on_product},
  their \emph{tensor product} is
  \[
    L \,\mathbin{\otimes}\, L' \;:=\;
      L \otimes_{\mathcal{O}_{C \times_S T}} L',
  \]
  the locally-trivial descent of the scheme-level tensor product
  (\ref{mk-lem:tensorobj_lift_onproduct}). The result again lies in the carrier
  because the tensor product of two locally-trivial modules is locally trivial
  (\ref{mk-lem:tensorobj_preserves_locally_trivial}). This is the operation that
  realises addition \([L] + [L'] := [L \otimes L']\) on isomorphism classes.
\end{definition}

\begin{lemma}[The structure-sheaf unit is locally trivial]\label{mk-lem:rel_pic_sharp_unit_loctriv}
  \dcref{custom}

  \uses{mk-def:line_bundle_on_product, mk-lem:pullback_unit_iso}
  The structure sheaf \(\mathcal{O}_{C \times_S T}\) (the monoidal unit) is
  locally trivial, hence defines an element of the carrier
  \(\mathrm{OnProduct}(\pi_C, \pi_T)\). Its isomorphism class is the zero
  element of the group of \ref{mk-lem:rel_pic_sharp_groupoid}.
\end{lemma}

\begin{proof}
  \uses{mk-lem:pullback_unit_iso}
 Local triviality is checked on an affine chart: for
  any point \(x\), choose an affine open \(W \ni x\); the restriction of the
  structure sheaf to \(W\) is its pullback along the inclusion
  \(W \hookrightarrow C \times_S T\), and the pullback of the structure sheaf is
  again the structure sheaf by the unconditional pullback--unit isomorphism
  (\ref{mk-lem:pullback_unit_iso}). Composing these two isomorphisms trivialises
  \((\mathcal{O}_{C \times_S T})|_W \cong \mathcal{O}_W\).
\end{proof}

\begin{lemma}[Uniqueness of the tensor inverse]\label{mk-lem:rel_pic_sharp_inverse_unique}
  \dcref{0B8K,custom}

  \uses{mk-lem:tensorobj_lift_onproduct, mk-lem:tensorobj_unit_iso,
    mk-lem:tensorobj_comm_iso, mk-lem:tensorobj_assoc_iso}
  \textit{Source: \cite{stacks-project}, Tag 0B8K (\S 17.25, invertible modules): any
  two \(\otimes\)-inverses of \(\mathcal{L}\) are each isomorphic to
  \(\SheafHom_{\mathcal{O}_X}(\mathcal{L}, \mathcal{O}_X)\), hence to one
  another.}
  Let \(M, N, N'\) be relative line bundles, and suppose there are
  isomorphisms
  \[
    M \otimes N \;\cong\; \mathcal{O}_{C \times_S T}
    \qquad\text{and}\qquad
    M \otimes N' \;\cong\; \mathcal{O}_{C \times_S T}.
  \]
  Then \(N \cong N'\). In other words, the tensor inverse of a relative line
  bundle is unique up to isomorphism.
\end{lemma}

\begin{proof}
  \uses{mk-lem:tensorobj_lift_onproduct, mk-lem:tensorobj_unit_iso,
    mk-lem:tensorobj_comm_iso, mk-lem:tensorobj_assoc_iso}
 One transports \(N\) across the chain of coherence
  isomorphisms
  \[
    N \;\cong\; N \otimes \mathcal{O}
      \;\cong\; N \otimes (M \otimes N')
      \;\cong\; (N \otimes M) \otimes N'
      \;\cong\; (M \otimes N) \otimes N'
      \;\cong\; \mathcal{O} \otimes N'
      \;\cong\; N',
  \]
  using the right unitor and left unitor (\ref{mk-lem:tensorobj_unit_iso}), the
  given trivialisations together with bifunctoriality of \(\otimes\)
  (\ref{mk-lem:tensorobj_lift_onproduct}), the associator
  (\ref{mk-lem:tensorobj_assoc_iso}), and the braiding
  (\ref{mk-lem:tensorobj_comm_iso}). This is the relative-line-bundle analogue of
  the standard uniqueness of inverses in the Picard group.
\end{proof}

\begin{definition}[Descended addition on the relative Picard quotient]\label{mk-def:rel_add}
  \dcref{custom}

  \uses{mk-def:rel_tensor_obj, mk-thm:relative_pic_quotient_well_defined,
    mk-lem:tensorobj_lift_onproduct}
  The addition on the relative Picard quotient
  \(\Pic^\sharp_{C/k}(T) = \Pic(C \times_S T)/\pi_T^*\Pic(T)\) is the descent of
  the tensor product \ref{mk-def:rel_tensor_obj} to isomorphism classes:
  \[
    [L] + [L'] \;:=\; [\,L \otimes L'\,].
  \]
\end{definition}

\begin{proof}
  \uses{mk-def:rel_tensor_obj, mk-lem:tensorobj_lift_onproduct}
 Well-definedness on the quotient
  (\ref{mk-thm:relative_pic_quotient_well_defined}) is bifunctoriality of the
  tensor product: replacing \(L, L'\) by isomorphic representatives
  \(L \cong M\), \(L' \cong M'\) yields an isomorphism
  \(L \otimes L' \cong M \otimes M'\) (\ref{mk-lem:tensorobj_lift_onproduct}),
  so the class \([L \otimes L']\) is independent of the chosen representatives.
\end{proof}

\begin{definition}[Descended negation on the relative Picard quotient]\label{mk-def:rel_neg}
  \dcref{custom}

  \uses{mk-def:rel_tensor_obj, mk-lem:rel_pic_sharp_inverse_unique,
    mk-lem:tensorobj_inverse_invertible}
  The negation on the relative Picard quotient is
  \[
    -[L] \;:=\; [\,L^{\mathrm{inv}}\,],
  \]
  where \(L^{\mathrm{inv}}\) is the locally-trivial tensor-inverse witness
  supplied by \ref{mk-lem:tensorobj_inverse_invertible}, together with its
  trivialisation \(L \otimes L^{\mathrm{inv}} \cong \mathcal{O}_{C \times_S T}\).
\end{definition}

\begin{proof}
  \uses{mk-lem:rel_pic_sharp_inverse_unique, mk-lem:tensorobj_inverse_invertible}
 Well-definedness on the quotient uses uniqueness of
  the tensor inverse (\ref{mk-lem:rel_pic_sharp_inverse_unique}): if
  \(L \cong M\), then the chosen inverse witnesses \(L^{\mathrm{inv}}\) and
  \(M^{\mathrm{inv}}\) of \ref{mk-lem:tensorobj_inverse_invertible} both trivialise
  the same class up to the isomorphism \(L \cong M\), hence
  \(L^{\mathrm{inv}} \cong M^{\mathrm{inv}}\), so \([L^{\mathrm{inv}}]\) is
  independent of the representative.
\end{proof}


\begin{definition}[The \(H_T\)-coset relation on locally-trivial bundles]\label{mk-lem:relpic_rel}
  \dcref{custom}

  \uses{mk-def:line_bundle_on_product, mk-def:pullback_along_projection}
  On the locally-trivial carrier \(\mathrm{OnProduct}(\pi_C, \pi_T)\) of
  \ref{mk-def:line_bundle_on_product}, define, for locally-trivial \(L, L'\) on
  \(C \times_k T\), the relation
  \[
    L \sim L' \quad\Longleftrightarrow\quad
      \exists\, N \text{ locally trivial on } T,\;
      L \;\cong\; \pi_T^* N \otimes L' ,
  \]
  where \(\pi_T^* N = \mathtt{pullbackAlongProjection}\,\pi_C\,\pi_T\,N\) is the
  pullback of \ref{mk-def:pullback_along_projection}. This \emph{multiplicative
  form} is the left-coset relation of \(H_T := \pi_T^*\Pic(T)\): it holds exactly
  when \([L]\) and \([L']\) differ by an element of \(H_T\) (equivalently
  \(L \otimes L'^{-1} \cong \pi_T^* N\)). It is coarser than the absolute
  iso-class relation \(\mathtt{RelPicPresheaf.preimage\_subgroup}\).
\end{definition}

\begin{lemma}[Iso-class inclusion into the \(H_T\)-coset relation]\label{mk-lem:relpic_rel_of_iso}
  \dcref{custom}

  \uses{mk-lem:relpic_rel, mk-def:pullback_along_projection, mk-lem:pullback_unit_iso, mk-lem:tensorobj_unit_iso}
  If \(L\) and \(L'\) have isomorphic underlying bundles, then \(L \sim L'\); in
  particular \(\sim\) is coarser than the absolute iso-class relation, which is
  what lets the objectwise coherence isomorphisms descend to the quotient.
\end{lemma}

\begin{proof}
  \uses{mk-lem:relpic_rel, mk-lem:pullback_unit_iso, mk-lem:tensorobj_unit_iso}
  Take \(N = \mathcal{O}_T\). Then
  \[
    \pi_T^* \mathcal{O}_T \otimes L'
      \cong \mathcal{O}_{C \times_k T} \otimes L'
      \cong L' \cong L ,
  \]
  the first isomorphism by the pullback--unit isomorphism
  \ref{mk-lem:pullback_unit_iso}, the second by the left unitor
  \ref{mk-lem:tensorobj_unit_iso}, and the last by hypothesis. Hence
  \(L \cong \pi_T^* \mathcal{O}_T \otimes L'\), so \(L \sim L'\).
\end{proof}

\begin{lemma}[Reflexivity of the \(H_T\)-relation]\label{mk-lem:relpic_setoid_refl}
  \dcref{custom}

  \uses{mk-lem:relpic_rel, mk-lem:relpic_rel_of_iso}
  For every locally-trivial \(L\) on \(C \times_k T\) one has \(L \sim L\).
\end{lemma}

\begin{proof}
  \uses{mk-lem:relpic_rel_of_iso}
  Immediate from \ref{mk-lem:relpic_rel_of_iso} applied to the identity
  isomorphism \(L \cong L\). Concretely, taking \(N = \mathcal{O}_T\) gives
  \(\pi_T^* \mathcal{O}_T \otimes L \cong \mathcal{O}_{C \times_k T} \otimes L
  \cong L\), so \(L \sim L\).
\end{proof}

\begin{lemma}[Symmetry of the \(H_T\)-relation]\label{mk-lem:relpic_setoid_symm}
  \dcref{custom}

  \uses{mk-lem:relpic_rel, mk-lem:tensorobj_inverse_invertible, mk-lem:tensorobj_assoc_iso, mk-lem:tensorobj_comm_iso, mk-lem:tensorobj_unit_iso, mk-lem:pullback_tensor_iso_loctriv, mk-lem:pullback_unit_iso}
  If \(L \sim L'\) then \(L' \sim L\).
\end{lemma}

\begin{proof}
  \uses{mk-lem:relpic_rel, mk-lem:tensorobj_inverse_invertible, mk-lem:tensorobj_assoc_iso, mk-lem:tensorobj_comm_iso, mk-lem:tensorobj_unit_iso, mk-lem:pullback_tensor_iso_loctriv, mk-lem:pullback_unit_iso}
  Suppose \(L \cong \pi_T^* N \otimes L'\) with \(N\) locally trivial, and let
  \(N^{-1}\) be its tensor inverse on \(T\)
  (\ref{mk-lem:tensorobj_inverse_invertible}), again locally trivial. Tensoring
  the hypothesis by \(\pi_T^*(N^{-1})\) and multiplying the pullbacks through
  the locally-trivial comparison isomorphism
  \ref{mk-lem:pullback_tensor_iso_loctriv},
  \[
    \pi_T^*(N^{-1}) \otimes L
      \cong \pi_T^*(N^{-1}) \otimes \pi_T^* N \otimes L'
      \cong \pi_T^*(N^{-1} \otimes N) \otimes L'
      \cong \pi_T^* \mathcal{O}_T \otimes L'
      \cong L' ,
  \]
  after reassociating and braiding with
  \ref{mk-lem:tensorobj_assoc_iso}, \ref{mk-lem:tensorobj_comm_iso}, and
  \ref{mk-lem:tensorobj_unit_iso}
  and applying the pullback--unit isomorphism \ref{mk-lem:pullback_unit_iso}.
  Hence \(L' \cong \pi_T^*(N^{-1}) \otimes L\) with \(N^{-1}\) locally trivial,
  so \(L' \sim L\).
\end{proof}

\begin{lemma}[Transitivity of the \(H_T\)-relation]\label{mk-lem:relpic_setoid_trans}
  \dcref{custom}

  \uses{mk-lem:relpic_rel, mk-lem:pullback_tensor_iso_loctriv, mk-lem:tensorobj_assoc_iso, mk-lem:tensorobj_preserves_locally_trivial}
  If \(L \sim L'\) and \(L' \sim L''\) then \(L \sim L''\).
\end{lemma}

\begin{proof}
  \uses{mk-lem:relpic_rel, mk-lem:pullback_tensor_iso_loctriv, mk-lem:tensorobj_assoc_iso, mk-lem:tensorobj_preserves_locally_trivial}
  Suppose \(L \cong \pi_T^* N \otimes L'\) and \(L' \cong \pi_T^* N' \otimes L''\)
  with \(N, N'\) locally trivial. Substituting and reassociating with the
  associator \ref{mk-lem:tensorobj_assoc_iso},
  \[
    L \cong \pi_T^* N \otimes (\pi_T^* N' \otimes L'')
      \cong (\pi_T^* N \otimes \pi_T^* N') \otimes L''
      \cong \pi_T^*(N \otimes N') \otimes L'' ,
  \]
  the last isomorphism by the locally-trivial comparison
  \ref{mk-lem:pullback_tensor_iso_loctriv}. Since \(N \otimes N'\) is locally
  trivial (\ref{mk-lem:tensorobj_preserves_locally_trivial}), \(L \sim L''\).
\end{proof}

\begin{lemma}[Middle-four interchange for \(\otimes\)]\label{mk-lem:tensor_middle_four}
  \dcref{custom}

  \uses{mk-lem:tensorobj_assoc_iso, mk-lem:tensorobj_comm_iso, mk-def:tensorobjisoofiso}
  For arbitrary modules \(A, B, C, D\) on a scheme \(X\) there is a canonical
  isomorphism
  \[
    (A \otimes B) \otimes (C \otimes D)
      \;\cong\;
    (A \otimes C) \otimes (B \otimes D)
  \]
  that interchanges the two inner factors. It is the coherence isomorphism used
  to show that the objectwise tensor product descends to a well-defined addition
  on the relative Picard quotient (\ref{mk-lem:relpic_add_welldef}).
\end{lemma}

\begin{proof}
  \uses{mk-lem:tensorobj_assoc_iso, mk-lem:tensorobj_comm_iso, mk-def:tensorobjisoofiso}
  The isomorphism is assembled entirely from the unconditional associator
  \ref{mk-lem:tensorobj_assoc_iso} and the braiding \ref{mk-lem:tensorobj_comm_iso},
  applied to a factor through the bifunctoriality of \(\otimes\)
  (\ref{mk-def:tensorobjisoofiso}); no coherence pentagon is invoked.
  Reassociating and braiding the middle pair,
  \[
    (A \otimes B) \otimes (C \otimes D)
      \cong A \otimes \bigl(B \otimes (C \otimes D)\bigr)
      \cong A \otimes \bigl((B \otimes C) \otimes D\bigr)
      \cong A \otimes \bigl((C \otimes B) \otimes D\bigr)
      \cong A \otimes \bigl(C \otimes (B \otimes D)\bigr)
      \cong (A \otimes C) \otimes (B \otimes D) ,
  \]
  where the first, third-to-last, and last isomorphisms are instances of the
  associator, the middle isomorphism braids \(B \otimes C \cong C \otimes B\)
  inside the fixed outer factors, and each step past the first is transported
  through \(\otimes\) by \ref{mk-def:tensorobjisoofiso}.
\end{proof}

\begin{lemma}[Addition is well defined on the \(H_T\)-quotient]\label{mk-lem:relpic_add_welldef}
  \dcref{custom}

  \uses{mk-lem:relpic_rel, mk-lem:tensor_middle_four, mk-lem:pullback_tensor_iso_loctriv, mk-lem:tensorobj_assoc_iso, mk-lem:tensorobj_comm_iso, mk-lem:tensorobj_preserves_locally_trivial}
  If \(L_1 \sim L_2\) and \(L_1' \sim L_2'\) then
  \(L_1 \otimes L_1' \sim L_2 \otimes L_2'\); hence the descended tensor product
  \(+\) is well defined on the quotient
  \(\mathrm{OnProduct}(\pi_C, \pi_T) / {\sim}\).
\end{lemma}

\begin{proof}
  \uses{mk-lem:relpic_rel, mk-lem:tensor_middle_four, mk-lem:pullback_tensor_iso_loctriv, mk-lem:tensorobj_assoc_iso, mk-lem:tensorobj_comm_iso, mk-lem:tensorobj_preserves_locally_trivial}
  Suppose \(L_1 \cong \pi_T^* N \otimes L_2\) and
  \(L_1' \cong \pi_T^* N' \otimes L_2'\) with \(N, N'\) locally trivial. By the
  middle-four interchange
  \((A \otimes B) \otimes (C \otimes D) \cong (A \otimes C) \otimes (B \otimes D)\)
  of \ref{mk-lem:tensor_middle_four}, assembled from the associator and the
  braiding \ref{mk-lem:tensorobj_assoc_iso} and \ref{mk-lem:tensorobj_comm_iso},
  \[
    L_1 \otimes L_1'
      \cong (\pi_T^* N \otimes L_2) \otimes (\pi_T^* N' \otimes L_2')
      \cong (\pi_T^* N \otimes \pi_T^* N') \otimes (L_2 \otimes L_2')
      \cong \pi_T^*(N \otimes N') \otimes (L_2 \otimes L_2') ,
  \]
  the last step by the locally-trivial comparison
  \ref{mk-lem:pullback_tensor_iso_loctriv}. As \(N \otimes N'\) is locally
  trivial (\ref{mk-lem:tensorobj_preserves_locally_trivial}),
  \(L_1 \otimes L_1' \sim L_2 \otimes L_2'\).
\end{proof}

\begin{definition}[The relative Picard carrier setoid]\label{mk-lem:relpic_setoid}
  \dcref{custom}

  \uses{mk-lem:relpic_rel, mk-lem:relpic_setoid_refl, mk-lem:relpic_setoid_symm, mk-lem:relpic_setoid_trans}
  The \(H_T\)-coset relation \(\sim\) of \ref{mk-lem:relpic_rel} is an equivalence
  relation on \(\mathrm{OnProduct}(\pi_C, \pi_T)\), by reflexivity
  \ref{mk-lem:relpic_setoid_refl}, symmetry \ref{mk-lem:relpic_setoid_symm}, and
  transitivity \ref{mk-lem:relpic_setoid_trans}. The resulting setoid presents the
  relative Picard quotient \(\Pic(C \times_k T) / H_T\), distinct from the
  absolute iso-class setoid \(\mathtt{RelPicPresheaf.preimage\_subgroup}\).
\end{definition}

\begin{lemma}[Negation is compatible with the \(H_T\)-relation]\label{mk-lem:relpic_rel_neg}
  \dcref{custom}

  \uses{mk-lem:relpic_rel, mk-lem:tensorobj_inverse_invertible, mk-lem:tensorobj_assoc_iso, mk-lem:tensorobj_comm_iso, mk-lem:tensorobj_unit_iso, mk-lem:pullback_tensor_iso_loctriv, mk-lem:pullback_unit_iso, mk-lem:rel_pic_sharp_inverse_unique}
  If \(L \sim M\), then their chosen tensor inverses satisfy \(L^{-1} \sim
  M^{-1}\). Hence the descended negation \([L] \mapsto [L^{-1}]\) is well defined
  on the quotient \(\mathrm{OnProduct}(\pi_C, \pi_T) / {\sim}\).
\end{lemma}

\begin{proof}
  \uses{mk-lem:relpic_rel, mk-lem:tensorobj_inverse_invertible, mk-lem:tensorobj_assoc_iso, mk-lem:tensorobj_comm_iso, mk-lem:tensorobj_unit_iso, mk-lem:pullback_tensor_iso_loctriv, mk-lem:pullback_unit_iso, mk-lem:rel_pic_sharp_inverse_unique}
  Suppose \(L \cong \pi_T^* N \otimes M\) with \(N\) locally trivial, and let
  \(N^{-1}\) be its tensor inverse (\ref{mk-lem:tensorobj_inverse_invertible}),
  again locally trivial. Write \(L^{-1}, M^{-1}\) for the chosen tensor inverses
  of \(L, M\). We show \(\pi_T^*(N^{-1}) \otimes M^{-1}\) is a tensor inverse of
  \(L\): tensoring the hypothesis with \(\pi_T^*(N^{-1}) \otimes M^{-1}\),
  reassociating and braiding with
  \ref{mk-lem:tensorobj_assoc_iso}, \ref{mk-lem:tensorobj_comm_iso}, and
  \ref{mk-lem:tensorobj_unit_iso},
  and multiplying the pullbacks through the locally-trivial comparison
  isomorphism \ref{mk-lem:pullback_tensor_iso_loctriv},
  \[
    L \otimes \bigl(\pi_T^*(N^{-1}) \otimes M^{-1}\bigr)
      \cong \pi_T^*(N \otimes N^{-1}) \otimes (M \otimes M^{-1})
      \cong \pi_T^* \mathcal{O}_T \otimes \mathcal{O}_{C \times_k T}
      \cong \mathcal{O}_{C \times_k T} ,
  \]
  using \(M \otimes M^{-1} \cong \mathcal{O}\) and the pullback--unit
  isomorphism \ref{mk-lem:pullback_unit_iso}. By uniqueness of tensor inverses
  (\ref{mk-lem:rel_pic_sharp_inverse_unique}),
  \(L^{-1} \cong \pi_T^*(N^{-1}) \otimes M^{-1}\), i.e.\ \(L^{-1} \sim M^{-1}\).
\end{proof}

\section{The relative Picard presheaf as a group-valued presheaf}
\label{mk-sec:relpic_groupvalued}

We now refine the set-valued functor of \ref{mk-thm:pullback_natural} into a
group-valued presheaf.

\begin{definition}[Relative Picard presheaf]\label{mk-def:rel_pic_sharp}
  \dcref{custom}

  \uses{mk-def:line_bundle_on_product, mk-def:pullback_along_projection,
    mk-thm:relative_pic_quotient_well_defined, mk-lem:relpic_setoid,
    mk-thm:rel_pic_addcommgroup_via_tensorobj, mk-lem:relpic_pullback_welldef,
    mk-thm:pullback_natural}
  \textit{Source: \cite{kleiman-picard}, ``The Picard scheme'', \S 2,
  Def.~df:Pfs.}
  Let \(C/k\) be a smooth proper geometrically integral curve. The
  \emph{relative Picard presheaf} of \(C\) over \(k\) is the functor
  \[
    \Pic^\sharp_{C/k} \;\;\colon\;\; (\Sch/k)^{op} \;\longrightarrow\; \Ab,
    \qquad
    T \;\longmapsto\; \Pic(C \times_k T) \,/\, \pi_T^*\Pic(T),
  \]
  where the right-hand side is the quotient abelian group of
  \ref{mk-thm:rel_pic_addcommgroup_via_tensorobj} on the \(H_T\)-coset
  carrier of \ref{mk-lem:relpic_setoid}, and the structure morphisms are the
  pullback maps \([L] \mapsto [g_C^* L]\), which descend to the quotient by
  \ref{mk-lem:relpic_pullback_welldef} and are group homomorphisms by
  \ref{mk-lem:rel_pic_sharp_functorial} below. In Kleiman's notation with
  \(X = C\) and \(S = \Spec k\), this is exactly \(\Pic_{X/S}\) of df:Pfs.
\end{definition}

\begin{lemma}[Naturality respects the group structure]\label{mk-lem:rel_pic_sharp_functorial}
  \dcref{custom}

  \uses{mk-def:rel_pic_sharp, mk-lem:pullback_compose,
    mk-lem:relpic_pullback_welldef, mk-lem:relpic_rel_of_iso,
    mk-lem:pullback_unit_iso, mk-lem:pullback_tensor_iso_loctriv,
    mk-thm:pullback_natural}
  \textit{Source: \cite{kleiman-picard}, ``The Picard scheme'', \S 2,
  Defs.~df:aPf + df:Pfs (the target category is the category of abelian
  groups; the structure maps are group homomorphisms).}
  For a morphism \(g : T' \to T\) of \(k\)-schemes, the assignment
  \([L] \mapsto [g_C^* L]\) is a well-defined homomorphism of abelian groups
  \(g^\sharp : \Pic^\sharp_{C/k}(T) \to \Pic^\sharp_{C/k}(T')\) with respect
  to the \(H_T\)-quotient structure of
  \ref{mk-thm:rel_pic_addcommgroup_via_tensorobj}, and the identity and
  composition laws \(\mathrm{id}^\sharp = \mathrm{id}\) and
  \((g \circ h)^\sharp = h^\sharp \circ g^\sharp\) hold as identities of group
  homomorphisms.
\end{lemma}
\begin{proof}


Well-definedness across \(H_T\)-cosets is
  \ref{mk-lem:relpic_pullback_welldef}. The map preserves zero because
  pullback preserves the structure sheaf (\ref{mk-lem:pullback_unit_iso})
  and preserves addition because pullback commutes with the tensor product
  of locally-trivial modules (\ref{mk-lem:pullback_tensor_iso_loctriv});
  both isomorphisms descend through \ref{mk-lem:relpic_rel_of_iso}. The
  identity law reduces to
  \(\mathrm{id}_C \times_k \mathrm{id}_T = \mathrm{id}\) and the invariance
  of classes under pullback along the identity; the composition law to
  \(\mathrm{id}_C \times_k (h \circ g) = (\mathrm{id}_C \times_k g) \circ
  (\mathrm{id}_C \times_k h)\) and pseudo-functoriality of pullback
  (\ref{mk-lem:pullback_compose}).
\end{proof}

\begin{theorem}[\(\Pic^\sharp_{C/k}\) is a group-valued presheaf]\label{mk-thm:rel_pic_sharp_presheaf}
  \dcref{custom}
  \uses{mk-def:rel_pic_sharp, mk-lem:rel_pic_sharp_functorial,
    mk-thm:rel_pic_addcommgroup_via_tensorobj}
  \textit{Source: \cite{kleiman-picard}, ``The Picard scheme'', \S 2,
  Defs.~df:aPf + df:Pfs.}
  The assignment \(T \mapsto \Pic^\sharp_{C/k}(T)\) of \ref{mk-def:rel_pic_sharp},
  together with the structure maps \(g \mapsto g^\sharp\) of
  \ref{mk-lem:rel_pic_sharp_functorial}, defines a presheaf of abelian groups
  \[
    \Pic^\sharp_{C/k} \;\;\colon\;\; (\Sch/k)^{op} \;\longrightarrow\;
                                     \AddCommGroup
  \]
  on the category of \(k\)-schemes.
\end{theorem}
\begin{proof}


The object map is well defined by
  \ref{mk-thm:rel_pic_addcommgroup_via_tensorobj}: for each \(k\)-scheme
  \(T\), the value \(\Pic^\sharp_{C/k}(T)\) is an abelian group. The
  morphism map \(g \mapsto g^\sharp\) is a group homomorphism by
  \ref{mk-lem:rel_pic_sharp_functorial}, and satisfies the identity and
  composition laws by the same lemma. These data assemble into a functor
  \((\Sch/k)^{op} \to \AddCommGroup\), i.e.\ a presheaf of abelian groups on
  \((\Sch/k)\).
\end{proof}

\subsection{Well-definedness, the absolute functor, and the quotient comparison}
\label{mk-subsec:relpic_relative_functor}

The well-definedness lemma below lets the pullback action descend to the
\(H_T\)-coset carrier of \ref{mk-lem:relpic_setoid}; it is the substrate of
the headline \ref{mk-lem:rel_pic_sharp_functorial}. The absolute iso-class
functor \(T \mapsto \Pic(C \times_k T)\) (Kleiman's \emph{absolute Picard
functor} df:aPf) is retained as \ref{mk-def:abs_pic_functor}, with the
natural componentwise quotient comparison
\ref{mk-lem:relpic_quotient_comparison} onto the relative functor of
\ref{mk-def:rel_pic_sharp}.

\begin{lemma}[Pullback descends to the \(H_T\)-quotient]\label{mk-lem:relpic_pullback_welldef}
  \dcref{custom}

  \uses{mk-lem:relpic_rel, mk-lem:pullback_compose,
    mk-lem:pullback_tensor_iso_loctriv, mk-lem:IsLocallyTrivial_pullback}
  \textit{Source: \cite{kleiman-picard}, ``The Picard scheme'', \S 2, Def.~df:Pfs
  (functoriality of the quotient).}
  Let \(g : T' \to T\) be a morphism of \(k\)-schemes and
  \(g_C := \mathrm{id}_C \times_k g\). If \(L \sim L'\) in the
  \(H_T\)-relation, then \(g_C^* L \sim g_C^* L'\) in the
  \(H_{T'}\)-relation.
\end{lemma}

\begin{proof}\uses{mk-lem:relpic_rel, mk-lem:pullback_compose, mk-lem:pullback_tensor_iso_loctriv}
  Let \(N\) be a coset witness, so \(L \cong \pi_T^* N \otimes L'\) with
  \(N\) locally trivial on \(T\). Applying \(g_C^*\) and the locally-trivial
  comparison isomorphism (\ref{mk-lem:pullback_tensor_iso_loctriv}),
  \(g_C^* L \cong g_C^* \pi_T^* N \otimes g_C^* L'\). By the pullback square
  \(\pi_T \circ g_C = g \circ \pi_{T'}\) (\ref{mk-lem:pullback_compose}),
  \(g_C^* \pi_T^* N \cong \pi_{T'}^* (g^* N)\), and \(g^* N\) is again
  locally trivial. Hence \(g^* N\) witnesses \(g_C^* L \sim g_C^* L'\).
\end{proof}

\begin{definition}[Absolute Picard functor]\label{mk-def:abs_pic_functor}
  \dcref{custom}

  \uses{mk-def:line_bundle_on_product, mk-lem:rel_pic_sharp_groupoid,
    mk-thm:pullback_natural}
  \textit{Source: \cite{kleiman-picard}, ``The Picard scheme'', \S 2, Def.~df:aPf
  (the absolute Picard functor).}
  The \emph{absolute Picard functor} of \(C\) over \(k\): the group-valued
  presheaf
  \[
    (\Sch/k)^{op} \longrightarrow \Ab, \qquad
    T \longmapsto \Pic(C \times_k T),
  \]
  on the iso-class carrier, with the tensor-product group law of
  \ref{mk-lem:rel_pic_sharp_groupoid} and the pullback structure maps of
  \ref{mk-thm:pullback_natural}. In Kleiman's notation this is \(\Pic_X\)
  of df:aPf (with \(X = C \times_k -\) over the test scheme).
  The object and morphism data are bundled compatibly with the pullback maps. It
  compares onto the relative functor of \ref{mk-def:rel_pic_sharp}
  via \ref{mk-lem:relpic_quotient_comparison}.
\end{definition}

\begin{lemma}[Quotient comparison from the absolute functor]\label{mk-lem:relpic_quotient_comparison}
  \dcref{custom}

  \uses{mk-def:abs_pic_functor, mk-def:rel_pic_sharp, mk-lem:relpic_rel_of_iso}
  \textit{Source: \cite{kleiman-picard}, ``The Picard scheme'', \S 2, Def.~df:Pfs
  (the quotient defining \(\Pic_{X/S}\)).}
  The componentwise coset-quotient maps
  \(\Pic(C \times_k T) \to \Pic(C \times_k T)/H_T\) assemble into a natural
  transformation from the absolute functor of \ref{mk-def:abs_pic_functor} to
  the relative functor of \ref{mk-def:rel_pic_sharp}, and each
  component is a surjective group homomorphism.
\end{lemma}

\begin{proof}\uses{mk-lem:relpic_rel_of_iso}
  Each component is the canonical quotient map, well defined because the
  \(H_T\)-relation is coarser than isomorphism
  (\ref{mk-lem:relpic_rel_of_iso}); it is a group homomorphism because both
  group laws descend the same tensor-product operations on representatives,
  and it is surjective as a quotient map. Naturality holds on
  representatives: both routes send \([L]\) to \([g_C^* L]\).
\end{proof}

\section{\'Etale sheafification}
\label{mk-sec:relpic_etale_sheaf}

The presheaf \(\Pic^\sharp_{C/k}\) is not \emph{a priori} a sheaf (Kleiman
\S 2, L1330), and over a general field it is not representable: Kleiman states
this outright at L5105--L5108 for the conic \(u^2+v^2+w^2=0\) in
\(\mathbb{P}^2_{\mathbb{R}}\) --- a smooth plane conic, hence smooth, proper and
geometrically integral, over a field where it has no rational point ---
because \(\Pic_{(X/\mathbb{R})\et}\) is representable by \S 4 while the two
functors differ (Exercise \texttt{ex:Pfs}).

Therefore, to obtain a representable functor in the sense of
\ref{mk-chap:Picard_FGAPicRepresentability}, one must replace
\(\Pic^\sharp_{C/k}\) by its \emph{associated \'etale sheaf}. Kleiman \S 2 names
this sheaf \(\Pic_{(X/S)\et}\) and singles it out as the functor that Kleiman's
main existence theorem in \S 4 represents.

\begin{definition}[\'Etale sheafification of the relative Picard presheaf]\label{mk-def:rel_pic_etale_sheafification}
  \dcref{custom}

  \uses{mk-def:abs_pic_functor}
  \textit{Source: \cite{kleiman-picard}, ``The Picard scheme'', \S 2,
  Def.~df:Pfs (the \'etale-sheaf notation \(\Pic_{(X/S)\et}\)) and the
  associated-sheaf paragraph immediately following df:Pfs.}
  Let \(C/k\) be a smooth proper geometrically integral curve. The
  \emph{\'etale-sheafified relative Picard functor} is the sheafification of
  the presheaf \(\Pic^\sharp_{C/k}\) (\ref{mk-thm:rel_pic_sharp_presheaf}) in the
  global \'etale topology on \((\Sch/k)\):
  \[
    \Pic^\sharp_{(C/k)\et} \;\;:=\;\;
       \bigl(\Pic^\sharp_{C/k}\bigr)^{\sim_{\et}}.
  \]
  Equivalently, \(\Pic^\sharp_{(C/k)\et}\) is the unique (up to canonical
  isomorphism) sheaf of abelian groups on the \'etale site of \(\Sch/k\)
  equipped with a presheaf morphism $\Pic^\sharp_{C/k} \to
  \Pic^\sharp_{(C/k)\et}$ which is universal among presheaf morphisms from
  \(\Pic^\sharp_{C/k}\) to \'etale sheaves of abelian groups. In Kleiman's
  notation this is \(\Pic_{(X/S)\et}\) with \(X = C\) and \(S = \Spec k\). When \(k\)
  is algebraically closed and \(T = \Spec k'\) for an algebraically closed field
  extension \(k'/k\), the natural map
  \(\Pic^\sharp_{C/k}(T) \to \Pic^\sharp_{(C/k)\et}(T)\) is a bijection (Kleiman
  Exercise ex:Alr, L1357--L1361).
\end{definition}

\begin{theorem}[\'Etale sheafification preserves the abelian-group structure]\label{mk-thm:rel_pic_etale_sheaf_group_structure}
  \dcref{custom}

  \uses{mk-def:rel_pic_etale_sheafification, mk-thm:rel_pic_sharp_presheaf,
    mk-thm:rel_pic_etale_sheaf_unit_canonical}
  \textit{Source: \cite{kleiman-picard}, ``The Picard scheme'', \S 2,
  Def.~df:Pfs together with the associated-sheaf paragraph immediately
  following (the \'etale sheaf \(\Pic_{(X/S)\et}\) is named together with the
  abelian-group-valued relative Picard functor, meaning the sheaf is taken in
  the category of abelian groups, and is universal among such sheaves
  receiving a presheaf morphism from \(\Pic_{X/S}\)).}
  This block records the bare existence of a morphism of presheaves of
  abelian groups
  \[
    \Pic^\sharp_{C/k} \;\;\longrightarrow\;\;
    \Pic^\sharp_{(C/k)\et}
  \]
  inside the category of presheaves
  \((\Sch/k)^{op} \to \AddCommGroup\); equivalently, the hom-set of
  presheaf morphisms from \(\Pic^\sharp_{C/k}\) to
  \(\Pic^\sharp_{(C/k)\et}\) is nonempty. The full canonical statement
  --- that this morphism is the canonical sheafification unit and
  carries the universal property characterising
  \(\Pic^\sharp_{(C/k)\et}\) as the initial \'etale sheaf of abelian
  groups receiving a presheaf morphism from \(\Pic^\sharp_{C/k}\), so
  in particular that the abelian-group structure on
  \(\Pic^\sharp_{(C/k)\et}(T)\) is uniquely determined --- is recorded
  separately as \ref{mk-thm:rel_pic_etale_sheaf_unit_canonical} below.
\end{theorem}
\begin{proof}


The canonical morphism from a presheaf to its associated sheaf in
  any Grothendieck topology (Stacks tag 00WI; the canonical
  sheafification unit) supplies, in particular, a morphism of
  presheaves \(\Pic^\sharp_{C/k} \to \Pic^\sharp_{(C/k)\et}\); the
  hom-set of presheaf morphisms between these two objects is therefore
  nonempty. The full universal property of this canonical morphism
  --- and consequently the claim that the abelian-group structure on
  the sheafification is uniquely pinned by it --- is the content of
  the forward-looking \ref{mk-thm:rel_pic_etale_sheaf_unit_canonical}
  below; see its proof sketch for the canonical formulation.
\end{proof}

\begin{theorem}[Canonical sheafification unit with universal property (forward-looking)]
  \label{mk-thm:rel_pic_etale_sheaf_unit_canonical}
  \dcref{custom}

  \uses{mk-def:rel_pic_sharp, mk-def:rel_pic_etale_sheafification}
  \textit{Source: \cite{kleiman-picard}, ``The Picard scheme'', \S 2,
  Def.~df:Pfs together with the associated-sheaf paragraph
  immediately following (the natural map from a presheaf to its
  associated sheaf, taken inside the category of abelian groups, is
  universal among presheaf morphisms from \(\Pic_{X/S}\) into \'etale
  sheaves of abelian groups).}
  Let \(C/k\) be a smooth proper geometrically integral curve. There
  exists a canonical morphism of presheaves of abelian groups
  \[
    \eta \;\colon\; \Pic^\sharp_{C/k} \;\longrightarrow\;
    \Pic^\sharp_{(C/k)\et}
  \]
  with the following universal property. For every \'etale sheaf
  \(\mathcal{F}\) of abelian groups on \((\Sch/k)\) and every morphism
  of presheaves of abelian groups
  \(\varphi : \Pic^\sharp_{C/k} \to \mathcal{F}\), there is a unique
  morphism of \'etale sheaves of abelian groups
  \(\tilde\varphi : \Pic^\sharp_{(C/k)\et} \to \mathcal{F}\) such that
  \(\tilde\varphi \circ \eta = \varphi\). Equivalently, the pair
  \(\bigl(\Pic^\sharp_{(C/k)\et}, \, \eta\bigr)\) is initial among
  pairs \((\mathcal{F}, \, \Pic^\sharp_{C/k} \to \mathcal{F})\)
  consisting of an \'etale sheaf of abelian groups on \((\Sch/k)\) and
  a morphism of presheaves of abelian groups from \(\Pic^\sharp_{C/k}\)
  to \(\mathcal{F}\). In particular, the abelian-group structure on
  \(\Pic^\sharp_{(C/k)\et}(T)\) at each \(k\)-scheme \(T\) is uniquely
  determined (up to canonical isomorphism) by the abelian-group
  structure on \(\Pic^\sharp_{C/k}(T)\) and the universal property of
  \(\eta\).
\end{theorem}
\begin{proof}

  In any Grothendieck topology, every presheaf has a canonical
  associated sheaf, and the canonical morphism from a presheaf to its
  associated sheaf is universal among morphisms from the presheaf into
  sheaves (Stacks tag 00WI; Kleiman \S 2, paragraph following
  df:Pfs). The \'etale site of \((\Sch/k)\) is a Grothendieck
  topology, so this applies in particular to \(\Pic^\sharp_{C/k}\) and
  yields the canonical set-valued sheafification
  \((\Pic^\sharp_{C/k})^{\sim_\et}\) together with the universal arrow
  out of \(\Pic^\sharp_{C/k}\). The forgetful functor
  \(\AddCommGroup \to \mathbf{Set}\) preserves and reflects the
  filtered colimits and finite limits used by the plus-construction,
  so performing the same plus-construction inside \(\AddCommGroup\)
  yields a sheaf of abelian groups whose underlying set-valued sheaf
  coincides with the set-valued sheafification; the abelian-group
  structure on the sheafification is transported back through this
  identification. The resulting morphism \(\eta\) is then by
  construction the universal arrow in the abelian-group-valued
  sheafification, which is the universal property in the form stated.
\end{proof}

The relationship with \ref{mk-chap:Picard_FGAPicRepresentability} is captured by
the slogan: the FGA Picard scheme of A.2.c represents
\(\Pic^\sharp_{(C/k)\et}\) (Kleiman \S 4, the main existence theorem: ``\dots represents
\(\Pic_{(X/S)\et}\)''), \emph{not} \(\Pic^\sharp_{C/k}\) directly. The \'etale
sheafification step of this chapter is precisely what makes the
representability theorem in A.2.c well-posed.
